\begin{register}{H}{SPI\_CMD\_REG}{0x0000}\label{regdesc:SPICMDREG}
\regfield{(reserved)}{7}{25}{0000000}%
\regfield{SPI\_USR}{1}{24}{{0}}%
\regfield{SPI\_UPDATE}{1}{23}{{0}}%
\regfield{(reserved)}{5}{18}{00000}%
\regfield{SPI\_CONF\_BITLEN}{18}{0}{{0}}%
\reglabel{Reset}\regnewline%
\begin{regdesc}\begin{reglist}
\label{fielddesc:SPICONFBITLEN}\item [SPI\_CONF\_BITLEN] 配置 SPI CONF 阶段的 SPI\_CLK 周期。 (R/W)\\
单位：SPI\_CLK 时钟周期。\\
可在 CONF 阶段配置。
\label{fielddesc:SPIUPDATE}\item [SPI\_UPDATE] 配置是否将 SPI 寄存器从 APB 时钟域同步到 SPI 模块时钟域。(WT)\\
\begin{itemize}
    \item 0：不同步
    \item 1：同步
\end{itemize}
该位仅用于 SPI 主机。 
\label{fielddesc:SPIUSR}\item [SPI\_USR] 配置是否使能用户自定义命令。(R/W/SC)\\
\begin{itemize}
    \item 0：不使能
    \item 1：使能
\end{itemize}
置位此位将触发一次 SPI 操作。操作结束后此位被自动清零。CONF\_buf 不可更改该配置。 
\end{reglist}\end{regdesc}
\end{register}


\begin{register}{H}{SPI\_ADDR\_REG}{0x0004}\label{regdesc:SPIADDRREG}
\regfield{SPI\_USR\_ADDR\_VALUE}{32}{0}{{0}}%
\reglabel{Reset}\regnewline%
\begin{regdesc}\begin{reglist}
\label{fielddesc:SPIUSRADDRVALUE}\item [SPI\_USR\_ADDR\_VALUE] 配置从机地址。(R/W)\\
可在 CONF 阶段配置。 
\end{reglist}\end{regdesc}
\end{register}


\begin{register}{H}{SPI\_USER\_REG}{0x0010}\label{regdesc:SPIUSERREG}
\regfield{SPI\_USR\_COMMAND}{1}{31}{{1}}%
\regfield{SPI\_USR\_ADDR}{1}{30}{{0}}%
\regfield{SPI\_USR\_DUMMY}{1}{29}{{0}}%
\regfield{SPI\_USR\_MISO}{1}{28}{{0}}%
\regfield{SPI\_USR\_MOSI}{1}{27}{{0}}%
\regfield{SPI\_USR\_DUMMY\_IDLE}{1}{26}{{0}}%
\regfield{SPI\_USR\_MOSI\_HIGHPART}{1}{25}{{0}}%
\regfield{SPI\_USR\_MISO\_HIGHPART}{1}{24}{{0}}%
\regfield{(reserved)}{6}{18}{000000}%
\regfield{SPI\_SIO}{1}{17}{{0}}%
\regfield{(reserved)}{1}{16}{0}%
\regfield{SPI\_USR\_CONF\_NXT}{1}{15}{{0}}%
\regfield{(reserved)}{1}{14}{0}%
\regfield{SPI\_FWRITE\_QUAD}{1}{13}{{0}}%
\regfield{SPI\_FWRITE\_DUAL}{1}{12}{{0}}%
\regfield{(reserved)}{2}{10}{00}%
\regfield{SPI\_CK\_OUT\_EDGE}{1}{9}{{0}}%
\regfield{SPI\_RSCK\_I\_EDGE}{1}{8}{{0}}%
\regfield{SPI\_CS\_SETUP}{1}{7}{{1}}%
\regfield{SPI\_CS\_HOLD}{1}{6}{{1}}%
\regfield{SPI\_TSCK\_I\_EDGE}{1}{5}{{0}}%
\regfield{(reserved)}{1}{4}{0}%
\regfield{SPI\_QPI\_MODE}{1}{3}{{0}}%
\regfield{(reserved)}{2}{1}{00}%
\regfield{SPI\_DOUTDIN}{1}{0}{{0}}%
\reglabel{Reset}\regnewline%
\begin{regdesc}\begin{reglist}
\label{fielddesc:SPIDOUTDIN}\item [SPI\_DOUTDIN] 配置是否使能全双工通信。(R/W)\\
\begin{itemize}
    \item 0：不使能
    \item 1：使能
\end{itemize}
可在 CONF 阶段配置。 
\label{fielddesc:SPIQPIMODE}\item [SPI\_QPI\_MODE] 配置是否使能 QPI 模式。(R/W/SS/SC)\\
\begin{itemize}
    \item 0：不使能
    \item 1：使能
\end{itemize}
SPI 主机和从机均支持该配置。可在 CONF 阶段配置。
\label{fielddesc:SPITSCKIEDGE}\item [SPI\_TSCK\_I\_EDGE] 配置 SPI 从机是否更改 TSCK 极性。(R/W)\\
\begin{itemize}
    \item 0: TSCK = SPI\_CK\_I
    \item 1: TSCK = !SPI\_CK\_I
\end{itemize}
\label{fielddesc:SPICSHOLD}\item [SPI\_CS\_HOLD] 配置在 SPI 处于完成 (DONE) 阶段时是否保持 SPI CS 拉低。(R/W)\\
\begin{itemize}
    \item 0：不保持拉低
    \item 1：保持拉低
\end{itemize}
可在 CONF 阶段配置。 
\label{fielddesc:SPICSSETUP}\item [SPI\_CS\_SETUP] 配置在 SPI 处于准备 (PREP) 阶段时是否使能 SPI CS。(R/W)\\
\begin{itemize}
    \item 0：不使能
    \item 1：使能
\end{itemize}
可在 CONF 阶段配置。
\label{fielddesc:SPIRSCKIEDGE}\item [SPI\_RSCK\_I\_EDGE] 配置 SPI 从机是否更改 RSCK 极性。(R/W)\\
\begin{itemize}
    \item 0: RSCK = !SPI\_CK\_I
    \item 1: RSCK = SPI\_CK\_I
\end{itemize}

\item [见下页]
\end{reglist}\end{regdesc}
\end{register}

\addtocounter{Regfloat}{-1}
\begin{register}{H}{SPI\_USER\_REG}{0x0010}
\begin{regdesc}\begin{reglist}
\item [接上页]
\label{fielddesc:SPICKOUTEDGE}\item [SPI\_CK\_OUT\_EDGE] 配置 SPI 时钟模式，与 \hyperref[fielddesc:SPICKIDLEEDGE]{SPI\_CK\_IDLE\_EDGE} 搭配使用。(R/W)\\ 
可在 CONF 阶段配置。更多信息见章节 \ref{GP-SPI2 master clock control}。
\label{fielddesc:SPIFWRITEDUAL}\item [SPI\_FWRITE\_DUAL] 配置在写操作 (DOUT) 阶段时读数据的方式是否为 2-bit 方式。(R/W)\\
\begin{itemize}
    \item 0：不使能
    \item 1：使能
\end{itemize}
可在 CONF 阶段配置。 
\label{fielddesc:SPIFWRITEQUAD}\item [SPI\_FWRITE\_QUAD] 配置在写操作 (DOUT) 阶段时读数据的方式是否为 4-bit 方式。(R/W)\\
\begin{itemize}
    \item 0：不使能
    \item 1：使能
\end{itemize}
可在 CONF 阶段配置。 
\label{fielddesc:SPIUSRCONFNXT}\item [SPI\_USR\_CONF\_NXT] 配置是否使能下一次传输事务的 CONF 阶段。(R/W)\\
\begin{itemize}
    \item 0：当前传输事务结束后，本次分段配置传输结束。或者，当前的传输模式不是分段配置传输。
    \item 1：本次分段配置传输继续进行，开始下一次传输事务。
\end{itemize}
可在 CONF 阶段配置。
\label{fielddesc:SPISIO}\item [SPI\_SIO] 配置是否使能 3 线半双工通信，其中 MOSI 和 MISO 信号共享一个管脚。(R/W)\\
\begin{itemize}
    \item 0：不使能
    \item 1：使能
\end{itemize}
可在 CONF 阶段配置。
\label{fielddesc:SPIUSRMISOHIGHPART}\item [SPI\_USR\_MISO\_HIGHPART] 配置在读数据阶段是否使能“高位模式”，即仅访问高位 buffer：\hyperref[regdesc:SPIW0REG]{SPI\_W8\_REG} \verb+~+ \hyperref[regdesc:SPIW0REG]{SPI\_W15\_REG}。(R/W)\\
\begin{itemize}
    \item 0：不使能
    \item 1：使能
\end{itemize}
可在 CONF 阶段配置。 

\label{fielddesc:SPIUSRMOSIHIGHPART}\item [SPI\_USR\_MOSI\_HIGHPART] 配置在写数据阶段是否使能“高位模式”，即仅访问高位 buffer：\hyperref[regdesc:SPIW0REG]{SPI\_W8\_REG} \verb+~+ \hyperref[regdesc:SPIW0REG]{SPI\_W15\_REG}。(R/W)\\
\begin{itemize}
    \item 0：不使能
    \item 1：使能
\end{itemize}
可在 CONF 阶段配置。 

\item [见下页]
\end{reglist}\end{regdesc}
\end{register}

\addtocounter{Regfloat}{-1}
\begin{register}{H}{SPI\_USER\_REG}{0x0010}
\begin{regdesc}\begin{reglist}
\item [接上页]


\label{fielddesc:SPIUSRDUMMYIDLE}\item [SPI\_USR\_DUMMY\_IDLE] 配置是否在 DUMMY 阶段禁用 SPI 时钟。 (R/W)\\
\begin{itemize}
    \item 0：不禁用
    \item 1：禁用
\end{itemize}
可在 CONF 阶段配置。
\label{fielddesc:SPIUSRMOSI}\item [SPI\_USR\_MOSI] 配置是否使能一次操作的写数据 (DOUT) 阶段。 (R/W)\\
\begin{itemize}
    \item 0：不使能
    \item 1：使能
\end{itemize}
可在 CONF 阶段配置。
\label{fielddesc:SPIUSRMISO}\item [SPI\_USR\_MISO] 配置是否使能一次操作的读数据 (DIN) 阶段。 (R/W)\\
\begin{itemize}
    \item 0：不使能
    \item 1：使能
\end{itemize}
可在 CONF 阶段配置。
\label{fielddesc:SPIUSRDUMMY}\item [SPI\_USR\_DUMMY] 配置是否使能一次操作的 DUMMY 阶段。 (R/W)\\
\begin{itemize}
    \item 0：不使能
    \item 1：使能
\end{itemize}
可在 CONF 阶段配置。
\label{fielddesc:SPIUSRADDR}\item [SPI\_USR\_ADDR] 配置是否使能一次操作的地址 (ADDR) 阶段。 (R/W)\\
\begin{itemize}
    \item 0：不使能
    \item 1：使能
\end{itemize}
\label{fielddesc:SPIUSRCOMMAND}\item [SPI\_USR\_COMMAND] 配置是否使能一次操作的命令 (CMD) 阶段。(R/W)\\
\begin{itemize}
    \item 0：不使能
    \item 1：使能
\end{itemize}
可在 CONF 阶段配置。
\end{reglist}\end{regdesc}
\end{register}


\begin{register}{H}{SPI\_USER1\_REG}{0x0014}\label{regdesc:SPIUSER1REG}
\regfield{SPI\_USR\_ADDR\_BITLEN}{5}{27}{{23}}%
\regfield{SPI\_CS\_HOLD\_TIME}{5}{22}{{0x1}}%
\regfield{SPI\_CS\_SETUP\_TIME}{5}{17}{{0}}%
\regfield{SPI\_MST\_WFULL\_ERR\_END\_EN}{1}{16}{{1}}%
\regfield{(reserved)}{8}{8}{00000000}%
\regfield{SPI\_USR\_DUMMY\_CYCLELEN}{8}{0}{{7}}%
\reglabel{Reset}\regnewline%
\begin{regdesc}\begin{reglist}
\label{fielddesc:SPIUSRDUMMYCYCLELEN}\item [SPI\_USR\_DUMMY\_CYCLELEN] 配置 DUMMY 阶段的时长。(R/W)\\
单位：SPI\_CLK 时钟周期。\\
此值为（预期周期数 - 1）。可在 CONF 阶段配置。
\label{fielddesc:SPIMSTWFULLERRENDEN}\item [SPI\_MST\_WFULL\_ERR\_END\_EN] 配置在主机全双工或半双工模式下，如果发生 SPI RX AFIFO 满错误，是否终止 SPI 传输。(R/W)\\
\begin{itemize}
    \item 0：不终止
    \item 1：终止
\end{itemize}
\label{fielddesc:SPICSSETUPTIME}\item [SPI\_CS\_SETUP\_TIME] 配置准备 (PREP) 阶段的时长。(R/W)\\
单位：SPI\_CLK 时钟周期。\\
此值等于预期周期数 - 1。此字段与 \hyperref[fielddesc:SPICSSETUP]{SPI\_CS\_SETUP} 搭配使用。可在 CONF 阶段配置。
\label{fielddesc:SPICSHOLDTIME}\item [SPI\_CS\_HOLD\_TIME] 配置 CS 管脚的延迟周期。(R/W)\\
单位：SPI\_CLK 时钟周期。\\
此字段与 \hyperref[fielddesc:SPICSHOLD]{SPI\_CS\_HOLD} 搭配使用。可在 CONF 阶段配置。
\label{fielddesc:SPIUSRADDRBITLEN}\item [SPI\_USR\_ADDR\_BITLEN] 配置地址阶段的位长。(R/W)\\
此值为（预期位数 - 1）。可在 CONF 阶段配置。 
\end{reglist}\end{regdesc}
\end{register}


\begin{register}{H}{SPI\_USER2\_REG}{0x0018}\label{regdesc:SPIUSER2REG}
\regfield{SPI\_USR\_COMMAND\_BITLEN}{4}{28}{{7}}%
\regfield{SPI\_MST\_REMPTY\_ERR\_END\_EN}{1}{27}{{1}}%
\regfield{(reserved)}{11}{16}{00000000000}%
\regfield{SPI\_USR\_COMMAND\_VALUE}{16}{0}{{0}}%
\reglabel{Reset}\regnewline%
\begin{regdesc}\begin{reglist}
\label{fielddesc:SPIUSRCOMMANDVALUE}\item [SPI\_USR\_COMMAND\_VALUE] 配置命令值。(R/W)\\
可在 CONF 阶段配置。 
\label{fielddesc:SPIMSTREMPTYERRENDEN}\item [SPI\_MST\_REMPTY\_ERR\_END\_EN] 配置在主机全双工或半双工模式下，如果发生 SPI TX AFIFO 空错误，是否终止 SPI 传输。(R/W)\\
\begin{itemize}
    \item 0：不终止
    \item 1：终止
\end{itemize}
\label{fielddesc:SPIUSRCOMMANDBITLEN}\item [SPI\_USR\_COMMAND\_BITLEN] 配置命令阶段的位长。(R/W)\\
此值为（预期位数 - 1）。可在 CONF 阶段配置。
\end{reglist}\end{regdesc}
\end{register}


\begin{register}{H}{SPI\_CTRL\_REG}{0x0008}\label{regdesc:SPICTRLREG}
\regfield{(reserved)}{5}{27}{00000}%
\regfield{SPI\_WR\_BIT\_ORDER}{2}{25}{{0}}%
\regfield{SPI\_RD\_BIT\_ORDER}{2}{23}{{0}}%
\regfield{(reserved)}{1}{22}{0}%
\regfield{SPI\_WP\_POL}{1}{21}{{1}}%
\regfield{SPI\_HOLD\_POL}{1}{20}{{1}}%
\regfield{SPI\_D\_POL}{1}{19}{{1}}%
\regfield{SPI\_Q\_POL}{1}{18}{{1}}%
\regfield{(reserved)}{2}{16}{00}%
\regfield{SPI\_FREAD\_QUAD}{1}{15}{{0}}%
\regfield{SPI\_FREAD\_DUAL}{1}{14}{{0}}%
\regfield{(reserved)}{4}{10}{0000}%
\regfield{SPI\_FCMD\_QUAD}{1}{9}{{0}}%
\regfield{SPI\_FCMD\_DUAL}{1}{8}{{0}}%
\regfield{(reserved)}{1}{7}{0}%
\regfield{SPI\_FADDR\_QUAD}{1}{6}{{0}}%
\regfield{SPI\_FADDR\_DUAL}{1}{5}{{0}}%
\regfield{(reserved)}{1}{4}{0}%
\regfield{SPI\_DUMMY\_OUT}{1}{3}{{0}}%
\regfield{(reserved)}{3}{0}{000}%
\reglabel{Reset}\regnewline%
\begin{regdesc}\begin{reglist}
\label{fielddesc:SPIDUMMYOUT}\item [SPI\_DUMMY\_OUT] 配置在 DUMMY 阶段是否输出 FSPI 总线信号。 (R/W)\\
\begin{itemize}
    \item 0：不输出
    \item 1：输出
\end{itemize}
可在 CONF 阶段配置。
\label{fielddesc:SPIFADDRDUAL}\item [SPI\_FADDR\_DUAL] 配置在地址 (ADDR) 阶段时是否采用 2-bit 模式。(R/W)\\
\begin{itemize}
    \item 0：不采用
    \item 1：采用
\end{itemize}
可在 CONF 阶段配置。
\label{fielddesc:SPIFADDRQUAD}\item [SPI\_FADDR\_QUAD] 配置在地址 (ADDR) 阶段时是否采用 4-bit 模式。(R/W)\\
\begin{itemize}
    \item 0：不采用
    \item 1：采用
\end{itemize}
可在 CONF 阶段配置。
\label{fielddesc:SPIFCMDDUAL}\item [SPI\_FCMD\_DUAL] 配置在命令 (CMD) 阶段时是否采用 2-bit 模式。(R/W)\\
\begin{itemize}
    \item 0：不采用
    \item 1：采用
\end{itemize}
可在 CONF 阶段配置。
\label{fielddesc:SPIFCMDQUAD}\item [SPI\_FCMD\_QUAD] 配置在命令 (CMD) 阶段是否采用 4-bit 模式。(R/W)\\
\begin{itemize}
    \item 0：不采用
    \item 1：采用
\end{itemize}
可在 CONF 阶段配置。
\label{fielddesc:SPIFREADDUAL}\item [SPI\_FREAD\_DUAL] 配置在读数据 (DIN) 阶段是否采用 2-bit 模式。(R/W)\\
\begin{itemize}
    \item 0：不采用
    \item 1：采用
\end{itemize}
可在 CONF 阶段配置。

\item [见下页]
\end{reglist}\end{regdesc}
\end{register}

\addtocounter{Regfloat}{-1}
\begin{register}{H}{SPI\_CTRL\_REG}{0x0008}
\begin{regdesc}\begin{reglist}
\item [接上页]

\label{fielddesc:SPIFREADQUAD}\item [SPI\_FREAD\_QUAD] 配置在读数据 (DIN) 阶段是否采用 4-bit 模式。(R/W)\\
\begin{itemize}
    \item 0：不采用
    \item 1：采用
\end{itemize}
可在 CONF 阶段配置。
\label{fielddesc:SPIQPOL}\item [SPI\_Q\_POL] 配置 MISO 的极性。(R/W)\\
\begin{itemize}
    \item 0：低
    \item 1：高
\end{itemize}
可在 CONF 阶段配置。
\label{fielddesc:SPIDPOL}\item [SPI\_D\_POL] 配置 MOSI 的极性。(R/W)\\
\begin{itemize}
    \item 0：低
    \item 1：高
\end{itemize}
可在 CONF 阶段配置。
\label{fielddesc:SPIHOLDPOL}\item [SPI\_HOLD\_POL] 配置在 SPI 空闲状态下 SPI\_HOLD 的输出值。(R/W)\\
\begin{itemize}
    \item 0：输出低电平
    \item 1：输出高电平
\end{itemize}
可在 CONF 阶段配置。
\label{fielddesc:SPIWPPOL}\item [SPI\_WP\_POL] 配置在 SPI 空闲状态下 WP 信号的输出值。(R/W)\\
\begin{itemize}
    \item 0：输出低电平
    \item 1：输出高电平
\end{itemize}
可在 CONF 阶段配置。
\label{fielddesc:SPIRDBITORDER}\item [SPI\_RD\_BIT\_ORDER] 配置读数据 (MISO) 阶段的比特顺序。(R/W)\\
\begin{itemize}
    \item 0：先读高有效位
    \item 1：先读低有效位
\end{itemize}
可在 CONF 阶段配置。
\label{fielddesc:SPIWRBITORDER}\item [SPI\_WR\_BIT\_ORDER] 配置命令 (CMD)、地址 (ADDR) 和写数据 (MOSI) 阶段的比特顺序。(R/W)\\
\begin{itemize}
    \item 0：先读高有效位
    \item 1：先读低有效位
\end{itemize}
可在 CONF 阶段配置。
\end{reglist}\end{regdesc}
\end{register}


\begin{register}{H}{SPI\_MS\_DLEN\_REG}{0x001C}\label{regdesc:SPIMSDLENREG}
\regfield{(reserved)}{14}{18}{00000000000000}%
\regfield{SPI\_MS\_DATA\_BITLEN}{18}{0}{{0}}%
\reglabel{Reset}\regnewline%
\begin{regdesc}\begin{reglist}
\label{fielddesc:SPIMSDATABITLEN}\item [SPI\_MS\_DATA\_BITLEN] 配置主机 DMA 控制或 CPU 控制的 SPI 传输的数据位长。或配置从机 DMA 控制的传输中接收数据的位长。(R/W)\\
该值等于需要的位长 - 1。可在 CONF 阶段配置。 
\end{reglist}\end{regdesc}
\end{register}


\begin{register}{H}{SPI\_MISC\_REG}{0x0020}\label{regdesc:SPIMISCREG}
\regfield{(reserved)}{1}{31}{{0}}%
\regfield{SPI\_CS\_KEEP\_ACTIVE}{1}{30}{{0}}%
\regfield{SPI\_CK\_IDLE\_EDGE}{1}{29}{{0}}%
\regfield{(reserved)}{5}{24}{00000}%
\regfield{SPI\_SLAVE\_CS\_POL}{1}{23}{{0}}%
\regfield{(reserved)}{10}{13}{0000000000}%
\regfield{SPI\_MASTER\_CS\_POL}{6}{7}{{0}}%
\regfield{SPI\_CK\_DIS}{1}{6}{{0}}%
\regfield{SPI\_CS5\_DIS}{1}{5}{{1}}%
\regfield{SPI\_CS4\_DIS}{1}{4}{{1}}%
\regfield{SPI\_CS3\_DIS}{1}{3}{{1}}%
\regfield{SPI\_CS2\_DIS}{1}{2}{{1}}%
\regfield{SPI\_CS1\_DIS}{1}{1}{{1}}%
\regfield{SPI\_CS0\_DIS}{1}{0}{{0}}%
\reglabel{Reset}\regnewline%
\begin{regdesc}\begin{reglist}
\label{fielddesc:SPICS0DIS}\item [SPI\_CS\regindex{n}\_DIS (\regindex{n} = 0 $\sim$ 5)] 配置是否禁用 SPI\_CS\regindex{n} 管脚。(R/W)\\
\begin{itemize}
    \item 0：SPI\_CS\regindex{n} 信号来自 CS\regindex{n} 管脚或输出至 CS\regindex{n} 管脚
    \item 1：禁用 CS\regindex{n}
\end{itemize}
可在 CONF 阶段配置。 
\label{fielddesc:SPICKDIS}\item [SPI\_CK\_DIS]配置是否禁用 SPI\_CK 输出信号。 (R/W)\\
0：使能 SPI\_CLK 输出信号。\\
1：停止 SPI\_CLK 输出信号；\\
可在 CONF 阶段配置。 
\label{fielddesc:SPIMASTERCSPOL}\item [SPI\_MASTER\_CS\_POL[\regindex{n}]] 配置主机 SPI\_CS\regindex{n} (\regindex{n} = 0 $\sim$ 5)。(R/W)\\
\begin{itemize}
    \item 0：SPI\_CS\regindex{n} 低电平有效
    \item 1：SPI\_CS\regindex{n} 高电平有效
\end{itemize} 
\label{fielddesc:SPISLAVECSPOL}\item [SPI\_SLAVE\_CS\_POL] 配置 SPI 从机输入信号 CS 的极性。(R/W)\\
\begin{itemize}
    \item 0：保持不变
    \item 1：反相
\end{itemize}
可在 CONF 阶段配置。
\item [见下页]
\end{reglist}\end{regdesc}
\end{register}

\addtocounter{Regfloat}{-1}
\begin{register}{H}{SPI\_MISC\_REG}{0x0020}
\begin{regdesc}\begin{reglist}
\item [接上页]
\label{fielddesc:SPICKIDLEEDGE}\item [SPI\_CK\_IDLE\_EDGE] 配置 SPI\_CLK 线在 GP-SPI2 空闲状态时是否保持高电平。(R/W)\\
\begin{itemize}
    \item 0：保持低电平
    \item 1：保持高电平
\end{itemize}
可在 CONF 阶段配置。 
\label{fielddesc:SPICSKEEPACTIVE}\item [SPI\_CS\_KEEP\_ACTIVE] 配置 SPI\_CS 是否保持低电平。(R/W)\\
\begin{itemize}
    \item 0：不保持低电平
    \item 1：保持低电平
\end{itemize}
可在 CONF 阶段配置。 
\end{reglist}\end{regdesc}
\end{register}


\begin{register}{H}{SPI\_DMA\_CONF\_REG}{0x0030}\label{regdesc:SPIDMACONFREG}
\regfield{SPI\_DMA\_AFIFO\_RST}{1}{31}{{0}}%
\regfield{SPI\_BUF\_AFIFO\_RST}{1}{30}{{0}}%
\regfield{SPI\_RX\_AFIFO\_RST}{1}{29}{{0}}%
\regfield{SPI\_DMA\_TX\_ENA}{1}{28}{{0}}%
\regfield{SPI\_DMA\_RX\_ENA}{1}{27}{{0}}%
\regfield{(reserved)}{5}{22}{00000}%
\regfield{SPI\_RX\_EOF\_EN}{1}{21}{{0}}%
\regfield{SPI\_SLV\_TX\_SEG\_TRANS\_CLR\_EN}{1}{20}{{0}}%
\regfield{SPI\_SLV\_RX\_SEG\_TRANS\_CLR\_EN}{1}{19}{{0}}%
\regfield{SPI\_DMA\_SLV\_SEG\_TRANS\_EN}{1}{18}{{0}}%
\regfield{(reserved)}{16}{2}{0000000000000000}%
\regfield{SPI\_DMA\_INFIFO\_FULL}{1}{1}{1}%
\regfield{SPI\_DMA\_OUTFIFO\_EMPTY}{1}{0}{1}%
\reglabel{Reset}\regnewline%
\begin{regdesc}\begin{reglist}
\label{fielddesc:SPIDMAOUTFIFOEMPTY}\item [SPI\_DMA\_OUTFIFO\_EMPTY] 表示 DMA TX FIFO 是否就绪。(RO)\\
\begin{itemize}
    \item 0：DMA TX FIFO 已就绪，可以发送数据
    \item 1：DMA TX FIFO 尚未就绪，不能发送数据
\end{itemize}
\label{fielddesc:SPIDMAINFIFOFULL}\item [SPI\_DMA\_INFIFO\_FULL] 表示 DMA RX FIFO 是否就绪。(RO)\\
\begin{itemize}
    \item 0：DMA RX FIFO 已就绪，可以接收数据
    \item 1：DMA RX FIFO 尚未就绪，不能接收数据
\end{itemize}
\label{fielddesc:SPIDMASLVSEGTRANSEN}\item [SPI\_DMA\_SLV\_SEG\_TRANS\_EN] 配置是否使能半双工通信方式下，DMA 控制的从机连续传输。(R/W)\\
\begin{itemize}
    \item 0：不使能
    \item 1：使能
\end{itemize}
\label{fielddesc:SPISLVRXSEGTRANSCLREN}\item [SPI\_SLV\_RX\_SEG\_TRANS\_CLR\_EN] 在从机连续传输中，如果 DMA RX buffer 小于实际接收的数据长度：(R/W)
\begin{itemize}
    \item 1：后续 Wr\_DMA 传输事务中传输的数据都不接收；
    \item 0：当前 Wr\_DMA 传输事务中传输的数据不接收，但在后续的 Wr\_DMA 传输事务中：
        \begin{itemize}
            \item 如果 DMA RX buffer 长度不为 0，则后续 Wr\_DMA 传输事务中传输的数据会被接收。
            \item 如果 DMA RX buffer 长度为 0，则后续 Wr\_DMA 传输事务中传输的数据不会被接收。
        \end{itemize}
\end{itemize}
\label{fielddesc:SPISLVTXSEGTRANSCLREN}\item [SPI\_SLV\_TX\_SEG\_TRANS\_CLR\_EN] 在从机连续传输中，如果 DMA TX buffer 小于实际发送的数据长度：(R/W)
\begin{itemize}
    \item 1：后续传输事务中传输的数据都不更新，即旧数据被重复发送；
    \item 0：当前传输事务中传输的数据不更新，但在后续的传输事务中：
            \begin{itemize}
            \item 如果有新的数据填充到 DMA TX FIFO，则将发送新数据。
            \item 如果没有新的数据填充到 DMA TX FIFO，则没有新数据被发送。
        \end{itemize}
\end{itemize}
\item [见下页]
\end{reglist}\end{regdesc}
\end{register}

\addtocounter{Regfloat}{-1}
\begin{register}{H}{SPI\_DMA\_CONF\_REG}{0x0030}
\begin{regdesc}\begin{reglist}
\item [接上页]
\label{fielddesc:SPIRXEOFEN}\item [SPI\_RX\_EOF\_EN] 1：在 DMA 控制的数据传输过程中，如果 DMA 传输的数据比特数等于 (\hyperref[fielddesc:SPIMSDATABITLEN]{SPI\_MS\_DATA\_BITLEN} + 1)，则硬件会置位 GDMA\_IN\_SUC\_EOF\_CH\textcolor{red}{\it{n}}\_INT\_RAW。0：在单次传输中，GDMA\_IN\_SUC\_EOF\_CH\textcolor{red}{\it{n}}\_INT\_RAW 由 \hyperref[int:SPITRANSDONEINT]{SPI\_TRANS\_DONE\_INT} 事件置位；或在分段配置传输模式下，由 \hyperref[int:SPIDMASEGTRANSDONEINT]{SPI\_DMA\_SEG\_TRANS\_DONE\_INT} 事件置位。 (R/W)
\label{fielddesc:SPIDMARXENA}\item [SPI\_DMA\_RX\_ENA] 配置是否使能 SPI DMA 控制的接收数据模式。 (R/W)\\
\begin{itemize}
    \item 0：不使能
    \item 1：使能
\end{itemize}
\label{fielddesc:SPIDMATXENA}\item [SPI\_DMA\_TX\_ENA] 配置是否使能 SPI DMA 控制的发送数据模式。(R/W) \\
\begin{itemize}
    \item 0：不使能
    \item 1：使能
\end{itemize}

\label{fielddesc:SPIRXAFIFORST}\item [SPI\_RX\_AFIFO\_RST] 配置是否复位图 \ref{fig:spi-master-data-flow-control} 和图 \ref{fig:spi-slave-data-flow-control} 中的 spi\_rx\_afifo。(WT)\\
\begin{itemize}
    \item 0：不复位
    \item 1：复位
\end{itemize}
spi\_rx\_afifo 将在 SPI 主机和从机传输中用于接收数据。 
\label{fielddesc:SPIBUFAFIFORST}\item [SPI\_BUF\_AFIFO\_RST] 配置是否复位图 \ref{fig:spi-master-data-flow-control} 和图 \ref{fig:spi-slave-data-flow-control} 中的 buf\_tx\_afifo。(WT)\\
\begin{itemize}
    \item 0：不复位
    \item 1：复位
\end{itemize}
buf\_tx\_afifo 将在 CPU 控制的从机传输或主机传输中用于发送数据。
\label{fielddesc:SPIDMAAFIFORST}\item [SPI\_DMA\_AFIFO\_RST] 配置是否复位图 \ref{fig:spi-master-data-flow-control} 和图 \ref{fig:spi-slave-data-flow-control} 中的 dma\_tx\_afifo。(WT)\\
\begin{itemize}
    \item 0：不复位
    \item 1：复位
\end{itemize}
dma\_tx\_afifo 在 DMA 控制的从机传输中用于发送数据。 
\end{reglist}\end{regdesc}
\end{register}


\begin{register}{H}{SPI\_SLAVE\_REG}{0x00E0}\label{regdesc:SPISLAVEREG}
\regfield{(reserved)}{2}{30}{00}%
\regfield{SPI\_MST\_FD\_WAIT\_DMA\_TX\_DATA}{1}{29}{0}%
\regfield{SPI\_USR\_CONF}{1}{28}{{0}}%
\regfield{SPI\_SOFT\_RESET}{1}{27}{{0}}%
\regfield{SPI\_SLAVE\_MODE}{1}{26}{{0}}%
\regfield{SPI\_DMA\_SEG\_MAGIC\_VALUE}{4}{22}{{10}}%
\regfield{(reserved)}{10}{12}{0000000000}%
\regfield{SPI\_SLV\_WRBUF\_BITLEN\_EN}{1}{11}{{0}}%
\regfield{SPI\_SLV\_RDBUF\_BITLEN\_EN}{1}{10}{{0}}%
\regfield{SPI\_SLV\_WRDMA\_BITLEN\_EN}{1}{9}{{0}}%
\regfield{SPI\_SLV\_RDDMA\_BITLEN\_EN}{1}{8}{{0}}%
\regfield{(reserved)}{4}{4}{0000}%
\regfield{SPI\_RSCK\_DATA\_OUT}{1}{3}{{0}}%
\regfield{SPI\_CLK\_MODE\_13}{1}{2}{{0}}%
\regfield{SPI\_CLK\_MODE}{2}{0}{{0}}%
\reglabel{Reset}\regnewline%
\begin{regdesc}\begin{reglist}
\label{fielddesc:SPICLKMODE}\item [SPI\_CLK\_MODE] 配置 SPI 时钟模式。 (R/W)
\begin{itemize}
    \item 0：CS 信号无效时，SPI 时钟关闭；
    \item 1：CS 信号无效后，SPI 时钟延迟一个时钟周期；
    \item 2：CS 信号无效后，SPI 时钟延迟两个时钟周期；
    \item 3：SPI 时钟一直有效。
\end{itemize}
可在 CONF 阶段配置。
\label{fielddesc:SPICLKMODE13}\item [SPI\_CLK\_MODE\_13] 配置时钟模式。 (R/W)
\begin{itemize}
    \item 0：支持 SPI 时钟模式 0 或 2，见表 \ref{table:clock_pol_slave}。 
    \item 1：支持 SPI 时钟模式 1 或 3，见表 \ref{table:clock_pol_slave}。 
\end{itemize}
\label{fielddesc:SPIRSCKDATAOUT}\item [SPI\_RSCK\_DATA\_OUT] 配置输出数据的时钟沿。(R/W)\\
\begin{itemize}
    \item 0：在 TSCK 上升沿输出数据
    \item 1：在 RSCK 上升沿输出数据
\end{itemize}
\label{fielddesc:SPISLVRDDMABITLENEN}\item [SPI\_SLV\_RDDMA\_BITLEN\_EN] 配置在 DMA 控制的 Rd\_DMA 传输过程中，是否使用 \hyperref[fielddesc:SPISLVDATABITLEN]{SPI\_SLV\_DATA\_BITLEN} 存储主机读取从机的数据位长。 (R/W)\\
\begin{itemize}
    \item 0：不使用
    \item 1：使用
\end{itemize}
\label{fielddesc:SPISLVWRDMABITLENEN}\item [SPI\_SLV\_WRDMA\_BITLEN\_EN] 配置在 DMA 控制的 Wr\_DMA 传输过程中，是否使用 \hyperref[fielddesc:SPISLVDATABITLEN]{SPI\_SLV\_DATA\_BITLEN} 存储主机向从机写数据的位长。 (R/W)\\
\begin{itemize}
    \item 0：不使用
    \item 1：使用
\end{itemize}
\item [见下页]
\end{reglist}\end{regdesc}
\end{register}

\addtocounter{Regfloat}{-1}
\begin{register}{H}{SPI\_SLAVE\_REG}{0x00E0}
\begin{regdesc}\begin{reglist}
\item [接上页]
\label{fielddesc:SPISLVRDBUFBITLENEN}\item [SPI\_SLV\_RDBUF\_BITLEN\_EN] 配置在 CPU 控制的 Rd\_BUF 传输过程中，是否使用 \hyperref[fielddesc:SPISLVDATABITLEN]{SPI\_SLV\_DATA\_BITLEN} 存储主机读取从机的数据位长。 (R/W)\\
\begin{itemize}
    \item 0：不使用
    \item 1：使用
\end{itemize}
\label{fielddesc:SPISLVWRBUFBITLENEN}\item [SPI\_SLV\_WRBUF\_BITLEN\_EN] 配置在 CPU 控制的 Wr\_BUF 传输过程中，是否使用 \hyperref[fielddesc:SPISLVDATABITLEN]{SPI\_SLV\_DATA\_BITLEN} 存储主机向从机写数据的位长。 (R/W)\\
\begin{itemize}
    \item 0：不使用
    \item 1：使用
\end{itemize}
\label{fielddesc:SPIDMASEGMAGICVALUE}\item [SPI\_DMA\_SEG\_MAGIC\_VALUE] 配置 DMA 控制的分段配置传输中位图表的 Magic 值。 (R/W)
\label{fielddesc:SPISLAVEMODE}\item [SPI\_SLAVE\_MODE] 配置 SPI 工作模式。(R/W)\\
\begin{itemize}
    \item 0：主机
    \item 1：从机
\end{itemize}
\label{fielddesc:SPISOFTRESET}\item [SPI\_SOFT\_RESET] 配置是否软件复位 SPI 时钟线、CS 线和数据线。(WT)\\
\begin{itemize}
    \item 0：不复位
    \item 1：复位
\end{itemize}
可在 CONF 阶段配置。 
\label{fielddesc:SPIUSRCONF}\item [SPI\_USR\_CONF] 配置是否使能当前 DMA 控制分段配置传输的 CONF 阶段。(R/W)\\
\begin{itemize}
    \item 0：无效，表明当前传输不是分段配置传输
    \item 1：使能，开始分段配置传输
\end{itemize}

\label{fielddesc:SPIMSTFDWAITDMATXDATA}\item [SPI\_MST\_FD\_WAIT\_DMA\_TX\_DATA] 配置在主机全双工模式下 GP-SPI2 是否先等待 DMA TX 数据准备好之后，再开始 SPI 传输。
\begin{itemize}
    \item 0：GP-SPI2 无需等待 DMA TX 数据即可开始 SPI 传输。
    \item 1：等待
\end{itemize}
\end{reglist}\end{regdesc}
\end{register}


\begin{register}{H}{SPI\_SLAVE1\_REG}{0x00E4}\label{regdesc:SPISLAVE1REG}
\regfield{SPI\_SLV\_LAST\_ADDR}{6}{26}{{0}}%
\regfield{SPI\_SLV\_LAST\_COMMAND}{8}{18}{{0}}%
\regfield{SPI\_SLV\_DATA\_BITLEN}{18}{0}{{0}}%
\reglabel{Reset}\regnewline%
\begin{regdesc}\begin{reglist}
\label{fielddesc:SPISLVDATABITLEN}\item [SPI\_SLV\_DATA\_BITLEN] 配置在 SPI 从机全双工和半双工传输中，传输的数据位长。 (R/W/SS)
\label{fielddesc:SPISLVLASTCOMMAND}\item [SPI\_SLV\_LAST\_COMMAND] 配置从机的命令值。 (R/W/SS)
\label{fielddesc:SPISLVLASTADDR}\item [SPI\_SLV\_LAST\_ADDR] 配置从机的地址值。 (R/W/SS)
\end{reglist}\end{regdesc}
\end{register}


\begin{register}{H}{SPI\_CLOCK\_REG}{0x000C}\label{regdesc:SPICLOCKREG}
\regfield{SPI\_CLK\_EQU\_SYSCLK}{1}{31}{{1}}%
\regfield{(reserved)}{9}{22}{000000000}%
\regfield{SPI\_CLKDIV\_PRE}{4}{18}{{0}}%
\regfield{SPI\_CLKCNT\_N}{6}{12}{{0x3}}%
\regfield{SPI\_CLKCNT\_H}{6}{6}{{0x1}}%
\regfield{SPI\_CLKCNT\_L}{6}{0}{{0x3}}%
\reglabel{Reset}\regnewline%
\begin{regdesc}\begin{reglist}
\label{fielddesc:SPICLKCNTL}\item [SPI\_CLKCNT\_L] 用作主机时，必须与 SPI\_CLKCNT\_N 相等。在从机模式下，必须为 0。可在 CONF 阶段配置。 (R/W)
\label{fielddesc:SPICLKCNTH}\item [SPI\_CLKCNT\_H] 配置用作主机时 SPI\_CLK（高电平）的占空比。(R/W)\\
建议将此值配置为 floor ((SPI\_CLKCNT\_N + 1)/2 - 1)。floor() 表示向下取整值，例如 floor(2.2) = 2。从机模式下，必须为 0。可在 CONF 阶段配置。 
\label{fielddesc:SPICLKCNTN}\item [SPI\_CLKCNT\_N] 配置用作主机时 SPI\_CLK 的分频系数。(R/W)\\
SPI\_CLK 频率为 $f_{\textrm{clk\_spi\_mst}}$/(SPI\_CLKDIV\_PRE + 1)/(SPI\_CLKCNT\_N + 1)。可在 CONF 阶段配置。 
\label{fielddesc:SPICLKDIVPRE}\item [SPI\_CLKDIV\_PRE] 配置用作主机时 SPI\_CLK 的预分频系数。(R/W)\\
可在 CONF 阶段配置。 
\label{fielddesc:SPICLKEQUSYSCLK}\item [SPI\_CLK\_EQU\_SYSCLK] 配置用作主机时 SPI\_CLK 频率是否与 clk\_spi\_mst 频率相同。(R/W)\\
\begin{itemize}
    \item 0：SPI\_CLK 为 clk\_spi\_mst 的分频时钟
    \item 1：SPI\_CLK 与 clk\_spi\_mst 频率相同
\end{itemize}
可在 CONF 阶段配置。 
\end{reglist}\end{regdesc}
\end{register}


\begin{register}{H}{SPI\_CLK\_GATE\_REG}{0x00E8}\label{regdesc:SPICLKGATEREG}
\regfield{(reserved)}{29}{3}{00000000000000000000000000000}%
\regfield{(reserved)}{1}{2}{{0}}%
\regfield{(reserved)}{1}{1}{{0}}%
\regfield{SPI\_CLK\_EN}{1}{0}{{0}}%
\reglabel{Reset}\regnewline%
\begin{regdesc}\begin{reglist}
\label{fielddesc:SPICLKEN}\item [SPI\_CLK\_EN] 配置是否使能时钟门控。 (R/W)
\begin{itemize}
    \item 0：不使能
    \item 1：使能
\end{itemize}
\end{reglist}\end{regdesc}
\end{register}

\begin{register}{H}{SPI\_DIN\_MODE\_REG}{0x0024}\label{regdesc:SPIDINMODEREG}
\regfield{(reserved)}{15}{17}{000000000000000}%
\regfield{SPI\_TIMING\_HCLK\_ACTIVE}{1}{16}{{0}}%
\regfield{(reserved)}{8}{8}{00000000}%
\regfield{SPI\_DIN3\_MODE}{2}{6}{{0}}%
\regfield{SPI\_DIN2\_MODE}{2}{4}{{0}}%
\regfield{SPI\_DIN1\_MODE}{2}{2}{{0}}%
\regfield{SPI\_DIN0\_MODE}{2}{0}{{0}}%
\reglabel{Reset}\regnewline%
\begin{regdesc}\begin{reglist}
\label{fielddesc:SPIDIN0MODE}\item [SPI\_DIN0\_MODE] 配置 FSPID 信号的输入模式。(R/W)
\begin{itemize}
    \item 0：无输入延迟
    \item 1：在 clk\_spi\_mst 的第 (\hyperref[fielddesc:SPIDIN0NUM]{SPI\_DIN0\_NUM} + 1) 个下降沿输入
    \item 2：在 clk\_hclk 的 (\hyperref[fielddesc:SPIDIN0NUM]{SPI\_DIN0\_NUM} + 1) 个上升沿延时加上一个 clk\_spi\_mst 上升沿延时的时刻输入
    \item 3：在 clk\_hclk 的 (\hyperref[fielddesc:SPIDIN0NUM]{SPI\_DIN0\_NUM} + 1) 个上升沿延时加上一个 clk\_spi\_mst 下降沿延时的时刻输入
\end{itemize}
可在 CONF 阶段配置。 
\label{fielddesc:SPIDIN1MODE}\item [SPI\_DIN1\_MODE] 配置 FSPIQ 信号的输入模式。 (R/W)
\begin{itemize}
    \item 0：无输入延迟
    \item 1：在 clk\_spi\_mst 的第 (\hyperref[fielddesc:SPIDIN1NUM]{SPI\_DIN1\_NUM} + 1) 个下降沿输入
    \item 2：在 clk\_hclk 的 (\hyperref[fielddesc:SPIDIN1NUM]{SPI\_DIN1\_NUM} + 1) 个上升沿延时加上一个 clk\_spi\_mst 上升沿延时的时刻输入
    \item 3：在 clk\_hclk 的 (\hyperref[fielddesc:SPIDIN1NUM]{SPI\_DIN1\_NUM} + 1) 个上升沿延时加上一个 clk\_spi\_mst 下降沿延时的时刻输入
\end{itemize}
可在 CONF 阶段配置。
\label{fielddesc:SPIDIN2MODE}\item [SPI\_DIN2\_MODE] 配置 FSPIWP 信号的输入模式。(R/W)
\begin{itemize}
    \item 0：无输入延迟
    \item 1：在 clk\_spi\_mst 的第 (\hyperref[fielddesc:SPIDIN2NUM]{SPI\_DIN2\_NUM} + 1) 个下降沿输入
    \item 2：在 clk\_hclk 的 (\hyperref[fielddesc:SPIDIN2NUM]{SPI\_DIN2\_NUM} + 1) 个上升沿延时加上一个 clk\_spi\_mst 上升沿延时的时刻输入
    \item 3：在 clk\_hclk 的 (\hyperref[fielddesc:SPIDIN2NUM]{SPI\_DIN2\_NUM} + 1) 个上升沿延时加上一个 clk\_spi\_mst 下降沿延时的时刻输入
\end{itemize}
可在 CONF 阶段配置。

\item [见下页]
\end{reglist}\end{regdesc}
\end{register}

\addtocounter{Regfloat}{-1}
\begin{register}{H}{SPI\_DIN\_MODE\_REG}{0x0024}
\begin{regdesc}\begin{reglist}
\item [接上页]
\label{fielddesc:SPIDIN3MODE}\item [SPI\_DIN3\_MODE] 配置 FSPIHD 信号的输入模式。 (R/W)
\begin{itemize}
    \item 0：无输入延迟
    \item 1：在 clk\_spi\_mst 的第 (\hyperref[fielddesc:SPIDIN3NUM]{SPI\_DIN3\_NUM}+1) 个下降沿输入
    \item 2：在 clk\_hclk 的 (\hyperref[fielddesc:SPIDIN3NUM]{SPI\_DIN3\_NUM}+1) 个上升沿延时加上一个 clk\_spi\_mst 上升沿延时的时刻输入
    \item 3：在 clk\_hclk 的 (\hyperref[fielddesc:SPIDIN3NUM]{SPI\_DIN3\_NUM}+1) 个上升沿延时加上一个 clk\_spi\_mst 下降沿延时的时刻输入
\end{itemize}
可在 CONF 阶段配置。

\label{fielddesc:SPITIMINGHCLKACTIVE}\item [SPI\_TIMING\_HCLK\_ACTIVE] 配置是否使能 SPI 输入信号时序模块的高频时钟 HCLK。(R/W)
\begin{itemize}
    \item 0：不使能
    \item 1：使能
\end{itemize}
可在 CONF 阶段配置。 
\end{reglist}\end{regdesc}
\end{register}


\begin{register}{H}{SPI\_DIN\_NUM\_REG}{0x0028}\label{regdesc:SPIDINNUMREG}
\regfield{(reserved)}{24}{8}{000000000000000000000000}%
\regfield{SPI\_DIN3\_NUM}{2}{6}{{0}}%
\regfield{SPI\_DIN2\_NUM}{2}{4}{{0}}%
\regfield{SPI\_DIN1\_NUM}{2}{2}{{0}}%
\regfield{SPI\_DIN0\_NUM}{2}{0}{{0}}%
\reglabel{Reset}\regnewline%
\begin{regdesc}\begin{reglist}
\label{fielddesc:SPIDIN0NUM}\item [SPI\_DIN0\_NUM] 配置输入信号 FSPID 的延迟周期数。具体的延迟模式由 \hyperref[fielddesc:SPIDIN0MODE]{SPI\_DIN0\_MODE} 配置。可在 CONF 阶段配置。 (R/W)
\begin{itemize}
    \item 0：延迟 1 个时钟周期
    \item 1：延迟 2 个时钟周期
    \item 2：延迟 3 个时钟周期
    \item 3：延迟 4 个时钟周期
\end{itemize}
\label{fielddesc:SPIDIN1NUM}\item [SPI\_DIN1\_NUM] 配置输入信号 FSPIQ 的延迟周期数。具体的延迟模式由 \hyperref[fielddesc:SPIDIN1MODE]{SPI\_DIN1\_MODE} 配置。可在 CONF 阶段配置。 (R/W)
\begin{itemize}
    \item 0：延迟 1 个时钟周期
    \item 1：延迟 2 个时钟周期
    \item 2：延迟 3 个时钟周期
    \item 3：延迟 4 个时钟周期
\end{itemize}
\label{fielddesc:SPIDIN2NUM}\item [SPI\_DIN2\_NUM] 配置输入信号 FSPIWP 的延迟周期数。具体的延迟模式由 \hyperref[fielddesc:SPIDIN2MODE]{SPI\_DIN2\_MODE} 配置。可在 CONF 阶段配置。 (R/W)
\begin{itemize}
    \item 0：延迟 1 个时钟周期
    \item 1：延迟 2 个时钟周期
    \item 2：延迟 3 个时钟周期
    \item 3：延迟 4 个时钟周期
\end{itemize}
\label{fielddesc:SPIDIN3NUM}\item [SPI\_DIN3\_NUM] 配置输入信号 FSPIHD 的延迟周期数。具体的延迟模式由 \hyperref[fielddesc:SPIDIN3MODE]{SPI\_DIN3\_MODE} 配置。可在 CONF 阶段配置。 (R/W)
\begin{itemize}
    \item 0：延迟 1 个时钟周期
    \item 1：延迟 2 个时钟周期
    \item 2：延迟 3 个时钟周期
    \item 3：延迟 4 个时钟周期
\end{itemize}
\end{reglist}\end{regdesc}
\end{register}


\begin{register}{H}{SPI\_DOUT\_MODE\_REG}{0x002C}\label{regdesc:SPIDOUTMODEREG}
\regfield{(reserved)}{28}{4}{0000000000000000000000000000}%
\regfield{SPI\_DOUT3\_MODE}{1}{3}{{0}}%
\regfield{SPI\_DOUT2\_MODE}{1}{2}{{0}}%
\regfield{SPI\_DOUT1\_MODE}{1}{1}{{0}}%
\regfield{SPI\_DOUT0\_MODE}{1}{0}{{0}}%
\reglabel{Reset}\regnewline%
\begin{regdesc}\begin{reglist}
\label{fielddesc:SPIDOUT0MODE}\item [SPI\_DOUT0\_MODE] 配置 FSPID 信号的输出模式。可在 CONF 阶段配置。 (R/W)
\begin{itemize}
    \item 0：无输出延迟
    \item 1：在 SPI 模块时钟的下降沿，延迟一个时钟周期后输出
\end{itemize}
\label{fielddesc:SPIDOUT1MODE}\item [SPI\_DOUT1\_MODE] 配置 FSPIQ 信号的输出模式。可在 CONF 阶段配置。 (R/W)
\begin{itemize}
    \item 0：无输出延迟
    \item 1：在 SPI 模块时钟的下降沿，延迟一个时钟周期后输出
\end{itemize}
\label{fielddesc:SPIDOUT2MODE}\item [SPI\_DOUT2\_MODE] 配置 FSPIWP 信号的输出模式。可在 CONF 阶段配置。 (R/W)
\begin{itemize}
    \item 0：无输出延迟
    \item 1：在 SPI 模块时钟的下降沿，延迟一个时钟周期后输出
\end{itemize}
\label{fielddesc:SPIDOUT3MODE}\item [SPI\_DOUT3\_MODE] 配置 FSPIHD 信号的输出模式。可在 CONF 阶段配置。 (R/W)
\begin{itemize}
    \item 0：无输出延迟
    \item 1：在 SPI 模块时钟的下降沿，延迟一个时钟周期后输出
\end{itemize}
\end{reglist}\end{regdesc}
\end{register}

\begin{register}{H}{SPI\_DMA\_INT\_ENA\_REG}{0x0034}\label{regdesc:SPIDMAINTENAREG}
\regfield{(reserved)}{11}{21}{00000000000}%
\regfield{SPI\_APP1\_INT\_ENA}{1}{20}{{0}}%
\regfield{SPI\_APP2\_INT\_ENA}{1}{19}{{0}}%
\regfield{SPI\_MST\_TX\_AFIFO\_REMPTY\_ERR\_INT\_ENA}{1}{18}{{0}}%
\regfield{SPI\_MST\_RX\_AFIFO\_WFULL\_ERR\_INT\_ENA}{1}{17}{{0}}%
\regfield{SPI\_SLV\_CMD\_ERR\_INT\_ENA}{1}{16}{{0}}%
\regfield{(reserved)}{1}{15}{{0}}%
\regfield{SPI\_SEG\_MAGIC\_ERR\_INT\_ENA}{1}{14}{{0}}%
\regfield{SPI\_DMA\_SEG\_TRANS\_DONE\_INT\_ENA}{1}{13}{{0}}%
\regfield{SPI\_TRANS\_DONE\_INT\_ENA}{1}{12}{{0}}%
\regfield{SPI\_SLV\_WR\_BUF\_DONE\_INT\_ENA}{1}{11}{{0}}%
\regfield{SPI\_SLV\_RD\_BUF\_DONE\_INT\_ENA}{1}{10}{{0}}%
\regfield{SPI\_SLV\_WR\_DMA\_DONE\_INT\_ENA}{1}{9}{{0}}%
\regfield{SPI\_SLV\_RD\_DMA\_DONE\_INT\_ENA}{1}{8}{{0}}%
\regfield{SPI\_SLV\_CMDA\_INT\_ENA}{1}{7}{{0}}%
\regfield{SPI\_SLV\_CMD9\_INT\_ENA}{1}{6}{{0}}%
\regfield{SPI\_SLV\_CMD8\_INT\_ENA}{1}{5}{{0}}%
\regfield{SPI\_SLV\_CMD7\_INT\_ENA}{1}{4}{{0}}%
\regfield{SPI\_SLV\_EN\_QPI\_INT\_ENA}{1}{3}{{0}}%
\regfield{SPI\_SLV\_EX\_QPI\_INT\_ENA}{1}{2}{{0}}%
\regfield{SPI\_DMA\_OUTFIFO\_EMPTY\_ERR\_INT\_ENA}{1}{1}{{0}}%
\regfield{SPI\_DMA\_INFIFO\_FULL\_ERR\_INT\_ENA}{1}{0}{{0}}%
\reglabel{Reset}\regnewline%
\begin{regdesc}\begin{reglist}
\label{fielddesc:SPIDMAINFIFOFULLERRINTENA}\item [SPI\_DMA\_INFIFO\_FULL\_ERR\_INT\_ENA] 写 1 使能 \hyperref[int:SPIINFIFOFULLERRINT]{SPI\_DMA\_INFIFO\_FULL\_ERR\_INT} 中断。 (R/W)
\label{fielddesc:SPIDMAOUTFIFOEMPTYERRINTENA}\item [SPI\_DMA\_OUTFIFO\_EMPTY\_ERR\_INT\_ENA] 写 1 使能 \hyperref[int:SPIOUTFIFOEMPTYERRINT]{SPI\_DMA\_OUTFIFO\_EMPTY\_ERR\_INT} 中断。 (R/W)
\label{fielddesc:SPISLVEXQPIINTENA}\item [SPI\_SLV\_EX\_QPI\_INT\_ENA] 写 1 使能 \hyperref[int:SPISLVEXQPIINT]{SPI\_SLV\_EX\_QPI\_INT} 中断。 (R/W)
\label{fielddesc:SPISLVENQPIINTENA}\item [SPI\_SLV\_EN\_QPI\_INT\_ENA] 写 1 使能 \hyperref[int:SPISLVENQPIINT]{SPI\_SLV\_EN\_QPI\_INT} 中断。 (R/W)
\label{fielddesc:SPISLVCMD7INTENA}\item [SPI\_SLV\_CMD7\_INT\_ENA] 写 1 使能 \hyperref[int:SPISLVCMD7INT]{SPI\_SLV\_CMD7\_INT} 中断。 (R/W)
\label{fielddesc:SPISLVCMD8INTENA}\item [SPI\_SLV\_CMD8\_INT\_ENA] 写 1 使能 \hyperref[int:SPISLVCMD8INT]{SPI\_SLV\_CMD8\_INT} 中断。 (R/W)
\label{fielddesc:SPISLVCMD9INTENA}\item [SPI\_SLV\_CMD9\_INT\_ENA] 写 1 使能 \hyperref[int:SPISLVCMD9INT]{SPI\_SLV\_CMD9\_INT} 中断。 (R/W)
\label{fielddesc:SPISLVCMDAINTENA}\item [SPI\_SLV\_CMDA\_INT\_ENA] 写 1 使能 \hyperref[int:SPISLVCMDAINT]{SPI\_SLV\_CMDA\_INT} 中断。 (R/W)
\label{fielddesc:SPISLVRDDMADONEINTENA}\item [SPI\_SLV\_RD\_DMA\_DONE\_INT\_ENA] 写 1 使能 \hyperref[int:SPISLVRDDMADONEINT]{SPI\_SLV\_RD\_DMA\_DONE\_INT} 中断。 (R/W)
\label{fielddesc:SPISLVWRDMADONEINTENA}\item [SPI\_SLV\_WR\_DMA\_DONE\_INT\_ENA] 写 1 使能 \hyperref[int:SPISLVWRDMADONEINT]{SPI\_SLV\_WR\_DMA\_DONE\_INT} 中断。 (R/W)
\label{fielddesc:SPISLVRDBUFDONEINTENA}\item [SPI\_SLV\_RD\_BUF\_DONE\_INT\_ENA] 写 1 使能 \hyperref[int:SPISLVRDBUFDONEINT]{SPI\_SLV\_RD\_BUF\_DONE\_INT} 中断。 (R/W)
\label{fielddesc:SPISLVWRBUFDONEINTENA}\item [SPI\_SLV\_WR\_BUF\_DONE\_INT\_ENA] 写 1 使能 \hyperref[int:SPISLVWRBUFDONEINT]{SPI\_SLV\_WR\_BUF\_DONE\_INT} 中断。 (R/W)
\label{fielddesc:SPITRANSDONEINTENA}\item [SPI\_TRANS\_DONE\_INT\_ENA] 写 1 使能 \hyperref[int:SPITRANSDONEINT]{SPI\_TRANS\_DONE\_INT} 中断。 (R/W)
\label{fielddesc:SPIDMASEGTRANSDONEINTENA}\item [SPI\_DMA\_SEG\_TRANS\_DONE\_INT\_ENA] 写 1 使能 \hyperref[int:SPIDMASEGTRANSDONEINT]{SPI\_DMA\_SEG\_TRANS\_DONE\_INT} 中断。 (R/W)
\label{fielddesc:SPISEGMAGICERRINTENA}\item [SPI\_SEG\_MAGIC\_ERR\_INT\_ENA] 写 1 使能 \hyperref[int:SPISEGMAGICERRINT]{SPI\_SEG\_MAGIC\_ERR\_INT} 中断。 (R/W)

\item [见下页]
\end{reglist}\end{regdesc}
\end{register}

\addtocounter{Regfloat}{-1}
\begin{register}{H}{SPI\_DMA\_INT\_ENA\_REG}{0x0034}
\begin{regdesc}\begin{reglist}
\item [接上页]

\label{fielddesc:SPISLVCMDERRINTENA}\item [SPI\_SLV\_CMD\_ERR\_INT\_ENA] 写 1 使能 \hyperref[int:SPISLVCMDERRINT]{SPI\_SLV\_CMD\_ERR\_INT} 中断。 (R/W)
\label{fielddesc:SPIMSTRXAFIFOWFULLERRINTENA}\item [SPI\_MST\_RX\_AFIFO\_WFULL\_ERR\_INT\_ENA] 写 1 使能 \hyperref[int:SPIMSTRXAFIFOWFULLERRINT]{SPI\_MST\_RX\_AFIFO\_WFULL\_ERR\_INT} 中断。 (R/W)
\label{fielddesc:SPIMSTTXAFIFOREMPTYERRINTENA}\item [SPI\_MST\_TX\_AFIFO\_REMPTY\_ERR\_INT\_ENA] 写 1 使能 \hyperref[int:SPIMSTTXAFIFOREMPTYERRINT]{SPI\_MST\_TX\_AFIFO\_REMPTY\_ERR\_INT} 中断。 (R/W)
\label{fielddesc:SPIAPP2INTENA}\item [SPI\_APP2\_INT\_ENA] 写 1 使能 \hyperref[int:SPIAPP2INT]{SPI\_APP2\_INT} 中断。 (R/W)
\label{fielddesc:SPIAPP1INTENA}\item [SPI\_APP1\_INT\_ENA] 写 1 使能 \hyperref[int:SPIAPP1INT]{SPI\_APP1\_INT} 中断。 (R/W)
\end{reglist}\end{regdesc}
\end{register}


\begin{register}{H}{SPI\_DMA\_INT\_CLR\_REG}{0x0038}\label{regdesc:SPIDMAINTCLRREG}
\regfield{(reserved)}{11}{21}{00000000000}%
\regfield{SPI\_APP1\_INT\_CLR}{1}{20}{{0}}%
\regfield{SPI\_APP2\_INT\_CLR}{1}{19}{{0}}%
\regfield{SPI\_MST\_TX\_AFIFO\_REMPTY\_ERR\_INT\_CLR}{1}{18}{{0}}%
\regfield{SPI\_MST\_RX\_AFIFO\_WFULL\_ERR\_INT\_CLR}{1}{17}{{0}}%
\regfield{SPI\_SLV\_CMD\_ERR\_INT\_CLR}{1}{16}{{0}}%
\regfield{(reserved)}{1}{15}{{0}}%
\regfield{SPI\_SEG\_MAGIC\_ERR\_INT\_CLR}{1}{14}{{0}}%
\regfield{SPI\_DMA\_SEG\_TRANS\_DONE\_INT\_CLR}{1}{13}{{0}}%
\regfield{SPI\_TRANS\_DONE\_INT\_CLR}{1}{12}{{0}}%
\regfield{SPI\_SLV\_WR\_BUF\_DONE\_INT\_CLR}{1}{11}{{0}}%
\regfield{SPI\_SLV\_RD\_BUF\_DONE\_INT\_CLR}{1}{10}{{0}}%
\regfield{SPI\_SLV\_WR\_DMA\_DONE\_INT\_CLR}{1}{9}{{0}}%
\regfield{SPI\_SLV\_RD\_DMA\_DONE\_INT\_CLR}{1}{8}{{0}}%
\regfield{SPI\_SLV\_CMDA\_INT\_CLR}{1}{7}{{0}}%
\regfield{SPI\_SLV\_CMD9\_INT\_CLR}{1}{6}{{0}}%
\regfield{SPI\_SLV\_CMD8\_INT\_CLR}{1}{5}{{0}}%
\regfield{SPI\_SLV\_CMD7\_INT\_CLR}{1}{4}{{0}}%
\regfield{SPI\_SLV\_EN\_QPI\_INT\_CLR}{1}{3}{{0}}%
\regfield{SPI\_SLV\_EX\_QPI\_INT\_CLR}{1}{2}{{0}}%
\regfield{SPI\_DMA\_OUTFIFO\_EMPTY\_ERR\_INT\_CLR}{1}{1}{{0}}%
\regfield{SPI\_DMA\_INFIFO\_FULL\_ERR\_INT\_CLR}{1}{0}{{0}}%
\reglabel{Reset}\regnewline%
\begin{regdesc}\begin{reglist}
\label{fielddesc:SPIDMAINFIFOFULLERRINTCLR}\item [SPI\_DMA\_INFIFO\_FULL\_ERR\_INT\_CLR] 写 1 清除 \hyperref[int:SPIINFIFOFULLERRINT]{SPI\_DMA\_INFIFO\_FULL\_ERR\_INT} 中断。 (WT)
\label{fielddesc:SPIDMAOUTFIFOEMPTYERRINTCLR}\item [SPI\_DMA\_OUTFIFO\_EMPTY\_ERR\_INT\_CLR] 写 1 清除 \hyperref[int:SPIOUTFIFOEMPTYERRINT]{SPI\_DMA\_OUTFIFO\_EMPTY\_ERR\_INT} 中断。 (WT)
\label{fielddesc:SPISLVEXQPIINTCLR}\item [SPI\_SLV\_EX\_QPI\_INT\_CLR] 写 1 清除 \hyperref[int:SPISLVEXQPIINT]{SPI\_SLV\_EX\_QPI\_INT} 中断。 (WT)
\label{fielddesc:SPISLVENQPIINTCLR}\item [SPI\_SLV\_EN\_QPI\_INT\_CLR] 写 1 清除 \hyperref[int:SPISLVENQPIINT]{SPI\_SLV\_EN\_QPI\_INT} 中断。 (WT)
\label{fielddesc:SPISLVCMD7INTCLR}\item [SPI\_SLV\_CMD7\_INT\_CLR] 写 1 清除 \hyperref[int:SPISLVCMD7INT]{SPI\_SLV\_CMD7\_INT} 中断。 (WT)
\label{fielddesc:SPISLVCMD8INTCLR}\item [SPI\_SLV\_CMD8\_INT\_CLR] 写 1 清除 \hyperref[int:SPISLVCMD8INT]{SPI\_SLV\_CMD8\_INT} 中断。 (WT)
\label{fielddesc:SPISLVCMD9INTCLR}\item [SPI\_SLV\_CMD9\_INT\_CLR] 写 1 清除 \hyperref[int:SPISLVCMD9INT]{SPI\_SLV\_CMD9\_INT} 中断。 (WT)
\label{fielddesc:SPISLVCMDAINTCLR}\item [SPI\_SLV\_CMDA\_INT\_CLR] 写 1 清除 \hyperref[int:SPISLVCMDAINT]{SPI\_SLV\_CMDA\_INT} 中断。 (WT)
\label{fielddesc:SPISLVRDDMADONEINTCLR}\item [SPI\_SLV\_RD\_DMA\_DONE\_INT\_CLR] 写 1 清除 \hyperref[int:SPISLVRDDMADONEINT]{SPI\_SLV\_RD\_DMA\_DONE\_INT} 中断。 (WT)
\label{fielddesc:SPISLVWRDMADONEINTCLR}\item [SPI\_SLV\_WR\_DMA\_DONE\_INT\_CLR] 写 1 清除 \hyperref[int:SPISLVWRDMADONEINT]{SPI\_SLV\_WR\_DMA\_DONE\_INT} 中断。 (WT)
\label{fielddesc:SPISLVRDBUFDONEINTCLR}\item [SPI\_SLV\_RD\_BUF\_DONE\_INT\_CLR] 写 1 清除 \hyperref[int:SPISLVRDBUFDONEINT]{SPI\_SLV\_RD\_BUF\_DONE\_INT} 中断。 (WT)
\label{fielddesc:SPISLVWRBUFDONEINTCLR}\item [SPI\_SLV\_WR\_BUF\_DONE\_INT\_CLR] 写 1 清除 \hyperref[int:SPISLVWRBUFDONEINT]{SPI\_SLV\_WR\_BUF\_DONE\_INT} 中断。 (WT)
\label{fielddesc:SPITRANSDONEINTCLR}\item [SPI\_TRANS\_DONE\_INT\_CLR] 写 1 清除 \hyperref[int:SPITRANSDONEINT]{SPI\_TRANS\_DONE\_INT} 中断。 (WT)
\label{fielddesc:SPIDMASEGTRANSDONEINTCLR}\item [SPI\_DMA\_SEG\_TRANS\_DONE\_INT\_CLR] 写 1 清除 \hyperref[int:SPIDMASEGTRANSDONEINT]{SPI\_DMA\_SEG\_TRANS\_DONE\_INT} 中断。 (WT)
\label{fielddesc:SPISEGMAGICERRINTCLR}\item [SPI\_SEG\_MAGIC\_ERR\_INT\_CLR] 写 1 清除 \hyperref[int:SPISEGMAGICERRINT]{SPI\_SEG\_MAGIC\_ERR\_INT} 中断。 (WT)

\item [见下页]
\end{reglist}\end{regdesc}
\end{register}

\addtocounter{Regfloat}{-1}
\begin{register}{H}{SPI\_DMA\_INT\_CLR\_REG}{0x0038}
\begin{regdesc}\begin{reglist}
\item [接上页]
\label{fielddesc:SPISLVCMDERRINTCLR}\item [SPI\_SLV\_CMD\_ERR\_INT\_CLR] 写 1 清除 \hyperref[int:SPISLVCMDERRINT]{SPI\_SLV\_CMD\_ERR\_INT} 中断。 (WT)
\label{fielddesc:SPIMSTRXAFIFOWFULLERRINTCLR}\item [SPI\_MST\_RX\_AFIFO\_WFULL\_ERR\_INT\_CLR] 写 1 清除 \hyperref[int:SPIMSTRXAFIFOWFULLERRINT]{SPI\_MST\_RX\_AFIFO\_WFULL\_ERR\_INT} 中断。 (WT)
\label{fielddesc:SPIMSTTXAFIFOREMPTYERRINTCLR}\item [SPI\_MST\_TX\_AFIFO\_REMPTY\_ERR\_INT\_CLR] 写 1 清除 \hyperref[int:SPIMSTTXAFIFOREMPTYERRINT]{SPI\_MST\_TX\_AFIFO\_REMPTY\_ERR\_INT} 中断。 (WT)
\label{fielddesc:SPIAPP2INTCLR}\item [SPI\_APP2\_INT\_CLR] 写 1 清除 \hyperref[int:SPIAPP2INT]{SPI\_APP2\_INT} 中断。 (WT)
\label{fielddesc:SPIAPP1INTCLR}\item [SPI\_APP1\_INT\_CLR] 写 1 清除 \hyperref[int:SPIAPP1INT]{SPI\_APP1\_INT} 中断。 (WT)
\end{reglist}\end{regdesc}
\end{register}


\begin{register}{H}{SPI\_DMA\_INT\_RAW\_REG}{0x003C}\label{regdesc:SPIDMAINTRAWREG}
\regfield{(reserved)}{11}{21}{00000000000}%
\regfield{SPI\_APP1\_INT\_RAW}{1}{20}{{0}}%
\regfield{SPI\_APP2\_INT\_RAW}{1}{19}{{0}}%
\regfield{SPI\_MST\_TX\_AFIFO\_REMPTY\_ERR\_INT\_RAW}{1}{18}{{0}}%
\regfield{SPI\_MST\_RX\_AFIFO\_WFULL\_ERR\_INT\_RAW}{1}{17}{{0}}%
\regfield{SPI\_SLV\_CMD\_ERR\_INT\_RAW}{1}{16}{{0}}%
\regfield{(reserved)}{1}{15}{{0}}%
\regfield{SPI\_SEG\_MAGIC\_ERR\_INT\_RAW}{1}{14}{{0}}%
\regfield{SPI\_DMA\_SEG\_TRANS\_DONE\_INT\_RAW}{1}{13}{{0}}%
\regfield{SPI\_TRANS\_DONE\_INT\_RAW}{1}{12}{{0}}%
\regfield{SPI\_SLV\_WR\_BUF\_DONE\_INT\_RAW}{1}{11}{{0}}%
\regfield{SPI\_SLV\_RD\_BUF\_DONE\_INT\_RAW}{1}{10}{{0}}%
\regfield{SPI\_SLV\_WR\_DMA\_DONE\_INT\_RAW}{1}{9}{{0}}%
\regfield{SPI\_SLV\_RD\_DMA\_DONE\_INT\_RAW}{1}{8}{{0}}%
\regfield{SPI\_SLV\_CMDA\_INT\_RAW}{1}{7}{{0}}%
\regfield{SPI\_SLV\_CMD9\_INT\_RAW}{1}{6}{{0}}%
\regfield{SPI\_SLV\_CMD8\_INT\_RAW}{1}{5}{{0}}%
\regfield{SPI\_SLV\_CMD7\_INT\_RAW}{1}{4}{{0}}%
\regfield{SPI\_SLV\_EN\_QPI\_INT\_RAW}{1}{3}{{0}}%
\regfield{SPI\_SLV\_EX\_QPI\_INT\_RAW}{1}{2}{{0}}%
\regfield{SPI\_DMA\_OUTFIFO\_EMPTY\_ERR\_INT\_RAW}{1}{1}{{0}}%
\regfield{SPI\_DMA\_INFIFO\_FULL\_ERR\_INT\_RAW}{1}{0}{{0}}%
\reglabel{Reset}\regnewline%
\begin{regdesc}\begin{reglist}
\label{fielddesc:SPIDMAINFIFOFULLERRINTRAW}\item [SPI\_DMA\_INFIFO\_FULL\_ERR\_INT\_RAW] \hyperref[int:SPIINFIFOFULLERRINT]{SPI\_DMA\_INFIFO\_FULL\_ERR\_INT} 的原始中断状态。 (R/WTC/SS)
\label{fielddesc:SPIDMAOUTFIFOEMPTYERRINTRAW}\item [SPI\_DMA\_OUTFIFO\_EMPTY\_ERR\_INT\_RAW] \hyperref[int:SPIOUTFIFOEMPTYERRINT]{SPI\_DMA\_OUTFIFO\_EMPTY\_ERR\_INT} 的原始中断状态。 (R/WTC/SS)
\label{fielddesc:SPISLVEXQPIINTRAW}\item [SPI\_SLV\_EX\_QPI\_INT\_RAW] \hyperref[int:SPISLVEXQPIINT]{SPI\_SLV\_EX\_QPI\_INT} 的原始中断状态。 (R/WTC/SS)
\label{fielddesc:SPISLVENQPIINTRAW}\item [SPI\_SLV\_EN\_QPI\_INT\_RAW] \hyperref[int:SPISLVENQPIINT]{SPI\_SLV\_EN\_QPI\_INT} 的原始中断状态。 (R/WTC/SS)
\label{fielddesc:SPISLVCMD7INTRAW}\item [SPI\_SLV\_CMD7\_INT\_RAW] \hyperref[int:SPISLVCMD7INT]{SPI\_SLV\_CMD7\_INT} 的原始中断状态。 (R/WTC/SS)
\label{fielddesc:SPISLVCMD8INTRAW}\item [SPI\_SLV\_CMD8\_INT\_RAW] \hyperref[int:SPISLVCMD8INT]{SPI\_SLV\_CMD8\_INT} 的原始中断状态。 (R/WTC/SS)
\label{fielddesc:SPISLVCMD9INTRAW}\item [SPI\_SLV\_CMD9\_INT\_RAW] \hyperref[int:SPISLVCMD9INT]{SPI\_SLV\_CMD9\_INT} 的原始中断状态。 (R/WTC/SS)
\label{fielddesc:SPISLVCMDAINTRAW}\item [SPI\_SLV\_CMDA\_INT\_RAW] \hyperref[int:SPISLVCMDAINT]{SPI\_SLV\_CMDA\_INT} 的原始中断状态。 (R/WTC/SS)
\label{fielddesc:SPISLVRDDMADONEINTRAW}\item [SPI\_SLV\_RD\_DMA\_DONE\_INT\_RAW] \hyperref[int:SPISLVRDDMADONEINT]{SPI\_SLV\_RD\_DMA\_DONE\_INT} 的原始中断状态。  (R/WTC/SS)
\label{fielddesc:SPISLVWRDMADONEINTRAW}\item [SPI\_SLV\_WR\_DMA\_DONE\_INT\_RAW] \hyperref[int:SPISLVWRDMADONEINT]{SPI\_SLV\_WR\_DMA\_DONE\_INT} 的原始中断状态。 (R/WTC/SS)
\label{fielddesc:SPISLVRDBUFDONEINTRAW}\item [SPI\_SLV\_RD\_BUF\_DONE\_INT\_RAW] \hyperref[int:SPISLVRDBUFDONEINT]{SPI\_SLV\_RD\_BUF\_DONE\_INT} 的原始中断状态。 (R/WTC/SS)
\label{fielddesc:SPISLVWRBUFDONEINTRAW}\item [SPI\_SLV\_WR\_BUF\_DONE\_INT\_RAW] \hyperref[int:SPISLVWRBUFDONEINT]{SPI\_SLV\_WR\_BUF\_DONE\_INT} 的原始中断状态。 (R/WTC/SS)
\label{fielddesc:SPITRANSDONEINTRAW}\item [SPI\_TRANS\_DONE\_INT\_RAW] \hyperref[int:SPITRANSDONEINT]{SPI\_TRANS\_DONE\_INT} 的原始中断状态。 (R/WTC/SS)
\item [见下页]
\end{reglist}\end{regdesc}
\end{register}

\addtocounter{Regfloat}{-1}
\begin{register}{H}{SPI\_DMA\_INT\_RAW\_REG}{0x003C}
\begin{regdesc}\begin{reglist}
\item [接上页]
\label{fielddesc:SPIDMASEGTRANSDONEINTRAW}\item [SPI\_DMA\_SEG\_TRANS\_DONE\_INT\_RAW] \hyperref[int:SPIDMASEGTRANSDONEINT]{SPI\_DMA\_SEG\_TRANS\_DONE\_INT} 的原始中断状态。 (R/WTC/SS)
\label{fielddesc:SPISEGMAGICERRINTRAW}\item [SPI\_SEG\_MAGIC\_ERR\_INT\_RAW] \hyperref[int:SPISEGMAGICERRINT]{SPI\_SEG\_MAGIC\_ERR\_INT} 的原始中断状态。 (R/WTC/SS)
\label{fielddesc:SPISLVCMDERRINTRAW}\item [SPI\_SLV\_CMD\_ERR\_INT\_RAW] \hyperref[int:SPISLVCMDERRINT]{SPI\_SLV\_CMD\_ERR\_INT} 的原始中断状态。 (R/WTC/SS)
\label{fielddesc:SPIMSTRXAFIFOWFULLERRINTRAW}\item [SPI\_MST\_RX\_AFIFO\_WFULL\_ERR\_INT\_RAW] \hyperref[int:SPIMSTRXAFIFOWFULLERRINT]{SPI\_MST\_RX\_AFIFO\_WFULL\_ERR\_INT} 的原始中断状态。 (R/WTC/SS)
\label{fielddesc:SPIMSTTXAFIFOREMPTYERRINTRAW}\item [SPI\_MST\_TX\_AFIFO\_REMPTY\_ERR\_INT\_RAW] \hyperref[int:SPIMSTTXAFIFOREMPTYERRINT]{SPI\_MST\_TX\_AFIFO\_REMPTY\_ERR\_INT} 的原始中断状态。 (R/WTC/SS)
\label{fielddesc:SPIAPP2INTRAW}\item [SPI\_APP2\_INT\_RAW] \hyperref[int:SPIAPP2INT]{SPI\_APP2\_INT} 的原始中断状态。该值仅由应用控制。 (R/WTC)
\label{fielddesc:SPIAPP1INTRAW}\item [SPI\_APP1\_INT\_RAW] \hyperref[int:SPIAPP1INT]{SPI\_APP1\_INT} 的原始中断状态。该值仅由应用控制。 (R/WTC)
\end{reglist}\end{regdesc}
\end{register}


\begin{register}{H}{SPI\_DMA\_INT\_ST\_REG}{0x0040}\label{regdesc:SPIDMAINTSTREG}
\regfield{(reserved)}{11}{21}{00000000000}%
\regfield{SPI\_APP1\_INT\_ST}{1}{20}{{0}}%
\regfield{SPI\_APP2\_INT\_ST}{1}{19}{{0}}%
\regfield{SPI\_MST\_TX\_AFIFO\_REMPTY\_ERR\_INT\_ST}{1}{18}{{0}}%
\regfield{SPI\_MST\_RX\_AFIFO\_WFULL\_ERR\_INT\_ST}{1}{17}{{0}}%
\regfield{SPI\_SLV\_CMD\_ERR\_INT\_ST}{1}{16}{{0}}%
\regfield{(reserved)}{1}{15}{{0}}%
\regfield{SPI\_SEG\_MAGIC\_ERR\_INT\_ST}{1}{14}{{0}}%
\regfield{SPI\_DMA\_SEG\_TRANS\_DONE\_INT\_ST}{1}{13}{{0}}%
\regfield{SPI\_TRANS\_DONE\_INT\_ST}{1}{12}{{0}}%
\regfield{SPI\_SLV\_WR\_BUF\_DONE\_INT\_ST}{1}{11}{{0}}%
\regfield{SPI\_SLV\_RD\_BUF\_DONE\_INT\_ST}{1}{10}{{0}}%
\regfield{SPI\_SLV\_WR\_DMA\_DONE\_INT\_ST}{1}{9}{{0}}%
\regfield{SPI\_SLV\_RD\_DMA\_DONE\_INT\_ST}{1}{8}{{0}}%
\regfield{SPI\_SLV\_CMDA\_INT\_ST}{1}{7}{{0}}%
\regfield{SPI\_SLV\_CMD9\_INT\_ST}{1}{6}{{0}}%
\regfield{SPI\_SLV\_CMD8\_INT\_ST}{1}{5}{{0}}%
\regfield{SPI\_SLV\_CMD7\_INT\_ST}{1}{4}{{0}}%
\regfield{SPI\_SLV\_EN\_QPI\_INT\_ST}{1}{3}{{0}}%
\regfield{SPI\_SLV\_EX\_QPI\_INT\_ST}{1}{2}{{0}}%
\regfield{SPI\_DMA\_OUTFIFO\_EMPTY\_ERR\_INT\_ST}{1}{1}{{0}}%
\regfield{SPI\_DMA\_INFIFO\_FULL\_ERR\_INT\_ST}{1}{0}{{0}}%
\reglabel{Reset}\regnewline%
\begin{regdesc}\begin{reglist}
\label{fielddesc:SPIDMAINFIFOFULLERRINTST}\item [SPI\_DMA\_INFIFO\_FULL\_ERR\_INT\_ST] \hyperref[int:SPIINFIFOFULLERRINT]{SPI\_DMA\_INFIFO\_FULL\_ERR\_INT} 的中断状态位。 (RO)
\label{fielddesc:SPIDMAOUTFIFOEMPTYERRINTST}\item [SPI\_DMA\_OUTFIFO\_EMPTY\_ERR\_INT\_ST] \hyperref[int:SPIOUTFIFOEMPTYERRINT]{SPI\_DMA\_OUTFIFO\_EMPTY\_ERR\_INT} 的中断状态位。 (RO)
\label{fielddesc:SPISLVEXQPIINTST}\item [SPI\_SLV\_EX\_QPI\_INT\_ST] \hyperref[int:SPISLVEXQPIINT]{SPI\_SLV\_EX\_QPI\_INT} 的中断状态位。 (RO)
\label{fielddesc:SPISLVENQPIINTST}\item [SPI\_SLV\_EN\_QPI\_INT\_ST] \hyperref[int:SPISLVENQPIINT]{SPI\_SLV\_EN\_QPI\_INT} 的中断状态位。 (RO)
\label{fielddesc:SPISLVCMD7INTST}\item [SPI\_SLV\_CMD7\_INT\_ST] \hyperref[int:SPISLVCMD7INT]{SPI\_SLV\_CMD7\_INT} 的中断状态位。 (RO)
\label{fielddesc:SPISLVCMD8INTST}\item [SPI\_SLV\_CMD8\_INT\_ST] \hyperref[int:SPISLVCMD8INT]{SPI\_SLV\_CMD8\_INT} 的中断状态位。 (RO)
\label{fielddesc:SPISLVCMD9INTST}\item [SPI\_SLV\_CMD9\_INT\_ST] \hyperref[int:SPISLVCMD9INT]{SPI\_SLV\_CMD9\_INT} 的中断状态位。 (RO)
\label{fielddesc:SPISLVCMDAINTST}\item [SPI\_SLV\_CMDA\_INT\_ST] \hyperref[int:SPISLVCMDAINT]{SPI\_SLV\_CMDA\_INT} 的中断状态位。 (RO)
\label{fielddesc:SPISLVRDDMADONEINTST}\item [SPI\_SLV\_RD\_DMA\_DONE\_INT\_ST] \hyperref[int:SPISLVRDDMADONEINT]{SPI\_SLV\_RD\_DMA\_DONE\_INT} 的中断状态位。 (RO)
\label{fielddesc:SPISLVWRDMADONEINTST}\item [SPI\_SLV\_WR\_DMA\_DONE\_INT\_ST] \hyperref[int:SPISLVWRDMADONEINT]{SPI\_SLV\_WR\_DMA\_DONE\_INT} 的中断状态位。 (RO)
\label{fielddesc:SPISLVRDBUFDONEINTST}\item [SPI\_SLV\_RD\_BUF\_DONE\_INT\_ST] \hyperref[int:SPISLVRDBUFDONEINT]{SPI\_SLV\_RD\_BUF\_DONE\_INT} 的中断状态位。 (RO)
\label{fielddesc:SPISLVWRBUFDONEINTST}\item [SPI\_SLV\_WR\_BUF\_DONE\_INT\_ST] \hyperref[int:SPISLVWRBUFDONEINT]{SPI\_SLV\_WR\_BUF\_DONE\_INT} 的中断状态位。 (RO)
\label{fielddesc:SPITRANSDONEINTST}\item [SPI\_TRANS\_DONE\_INT\_ST] \hyperref[int:SPITRANSDONEINT]{SPI\_TRANS\_DONE\_INT} 的中断状态位。 (RO)
\label{fielddesc:SPIDMASEGTRANSDONEINTST}\item [SPI\_DMA\_SEG\_TRANS\_DONE\_INT\_ST] \hyperref[int:SPIDMASEGTRANSDONEINT]{SPI\_DMA\_SEG\_TRANS\_DONE\_INT} 的中断状态位。 (RO)
\label{fielddesc:SPISEGMAGICERRINTST}\item [SPI\_SEG\_MAGIC\_ERR\_INT\_ST] \hyperref[int:SPISEGMAGICERRINT]{SPI\_SEG\_MAGIC\_ERR\_INT} 的中断状态位。 (RO)
\label{fielddesc:SPISLVCMDERRINTST}\item [SPI\_SLV\_CMD\_ERR\_INT\_ST] \hyperref[int:SPISLVCMDERRINT]{SPI\_SLV\_CMD\_ERR\_INT} 的中断状态位。 (RO)
\item [见下页]
\end{reglist}\end{regdesc}
\end{register}

\addtocounter{Regfloat}{-1}
\begin{register}{H}{SPI\_DMA\_INT\_ST\_REG}{0x0040}
\begin{regdesc}\begin{reglist}
\item [接上页]
\label{fielddesc:SPIMSTRXAFIFOWFULLERRINTST}\item [SPI\_MST\_RX\_AFIFO\_WFULL\_ERR\_INT\_ST] \hyperref[int:SPIMSTRXAFIFOWFULLERRINT]{SPI\_MST\_RX\_AFIFO\_WFULL\_ERR\_INT} 的中断状态位。 (RO)
\label{fielddesc:SPIMSTTXAFIFOREMPTYERRINTST}\item [SPI\_MST\_TX\_AFIFO\_REMPTY\_ERR\_INT\_ST] \hyperref[int:SPIMSTTXAFIFOREMPTYERRINT]{SPI\_MST\_TX\_AFIFO\_REMPTY\_ERR\_INT} 的中断状态位。 (RO)
\label{fielddesc:SPIAPP2INTST}\item [SPI\_APP2\_INT\_ST] \hyperref[int:SPIAPP2INT]{SPI\_APP2\_INT} 的中断状态位。 (RO)
\label{fielddesc:SPIAPP1INTST}\item [SPI\_APP1\_INT\_ST] \hyperref[int:SPIAPP1INT]{SPI\_APP1\_INT} 的中断状态位。 (RO)
\end{reglist}\end{regdesc}
\end{register}

\begin{register}{H}{SPI\_DMA\_INT\_SET\_REG}{0x0044}\label{regdesc:SPIDMAINTSETREG}
\regfield{(reserved)}{11}{21}{00000000000}%
\regfield{SPI\_APP1\_INT\_SET}{1}{20}{{0}}%
\regfield{SPI\_APP2\_INT\_SET}{1}{19}{{0}}%
\regfield{SPI\_MST\_TX\_AFIFO\_REMPTY\_ERR\_INT\_SET}{1}{18}{{0}}%
\regfield{SPI\_MST\_RX\_AFIFO\_WFULL\_ERR\_INT\_SET}{1}{17}{{0}}%
\regfield{SPI\_SLV\_CMD\_ERR\_INT\_SET}{1}{16}{{0}}%
\regfield{(reserved)}{1}{15}{{0}}%
\regfield{SPI\_SEG\_MAGIC\_ERR\_INT\_SET}{1}{14}{{0}}%
\regfield{SPI\_DMA\_SEG\_TRANS\_DONE\_INT\_SET}{1}{13}{{0}}%
\regfield{SPI\_TRANS\_DONE\_INT\_SET}{1}{12}{{0}}%
\regfield{SPI\_SLV\_WR\_BUF\_DONE\_INT\_SET}{1}{11}{{0}}%
\regfield{SPI\_SLV\_RD\_BUF\_DONE\_INT\_SET}{1}{10}{{0}}%
\regfield{SPI\_SLV\_WR\_DMA\_DONE\_INT\_SET}{1}{9}{{0}}%
\regfield{SPI\_SLV\_RD\_DMA\_DONE\_INT\_SET}{1}{8}{{0}}%
\regfield{SPI\_SLV\_CMDA\_INT\_SET}{1}{7}{{0}}%
\regfield{SPI\_SLV\_CMD9\_INT\_SET}{1}{6}{{0}}%
\regfield{SPI\_SLV\_CMD8\_INT\_SET}{1}{5}{{0}}%
\regfield{SPI\_SLV\_CMD7\_INT\_SET}{1}{4}{{0}}%
\regfield{SPI\_SLV\_EN\_QPI\_INT\_SET}{1}{3}{{0}}%
\regfield{SPI\_SLV\_EX\_QPI\_INT\_SET}{1}{2}{{0}}%
\regfield{SPI\_DMA\_OUTFIFO\_EMPTY\_ERR\_INT\_SET}{1}{1}{{0}}%
\regfield{SPI\_DMA\_INFIFO\_FULL\_ERR\_INT\_SET}{1}{0}{{0}}%
\reglabel{Reset}\regnewline%
\begin{regdesc}\begin{reglist}
\label{fielddesc:SPIDMAINFIFOFULLERRINTSET}\item [SPI\_DMA\_INFIFO\_FULL\_ERR\_INT\_SET] 写 1 置位 \hyperref[int:SPIINFIFOFULLERRINT]{SPI\_DMA\_INFIFO\_FULL\_ERR\_INT} 中断。 (WT)
\label{fielddesc:SPIDMAOUTFIFOEMPTYERRINTSET}\item [SPI\_DMA\_OUTFIFO\_EMPTY\_ERR\_INT\_SET] 写 1 置位 \hyperref[int:SPIOUTFIFOEMPTYERRINT]{SPI\_DMA\_OUTFIFO\_EMPTY\_ERR\_INT} 中断。 (WT)
\label{fielddesc:SPISLVEXQPIINTSET}\item [SPI\_SLV\_EX\_QPI\_INT\_SET] 写 1 置位 \hyperref[int:SPISLVEXQPIINT]{SPI\_SLV\_EX\_QPI\_INT} 中断。 (WT)
\label{fielddesc:SPISLVENQPIINTSET}\item [SPI\_SLV\_EN\_QPI\_INT\_SET] 写 1 置位 \hyperref[int:SPISLVENQPIINT]{SPI\_SLV\_EN\_QPI\_INT} 中断。 (WT)
\label{fielddesc:SPISLVCMD7INTSET}\item [SPI\_SLV\_CMD7\_INT\_SET] 写 1 置位 \hyperref[int:SPISLVCMD7INT]{SPI\_SLV\_CMD7\_INT} 中断。 (WT)
\label{fielddesc:SPISLVCMD8INTSET}\item [SPI\_SLV\_CMD8\_INT\_SET] 写 1 置位 \hyperref[int:SPISLVCMD8INT]{SPI\_SLV\_CMD8\_INT} 中断。 (WT)
\label{fielddesc:SPISLVCMD9INTSET}\item [SPI\_SLV\_CMD9\_INT\_SET] 写 1 置位 \hyperref[int:SPISLVCMD9INT]{SPI\_SLV\_CMD9\_INT} 中断。 (WT)
\label{fielddesc:SPISLVCMDAINTSET}\item [SPI\_SLV\_CMDA\_INT\_SET] 写 1 置位 \hyperref[int:SPISLVCMDAINT]{SPI\_SLV\_CMDA\_INT} 中断。 (WT)

\item [见下页]
\end{reglist}\end{regdesc}
\end{register}

\addtocounter{Regfloat}{-1}
\begin{register}{H}{SPI\_DMA\_INT\_SET\_REG}{0x0044}
\begin{regdesc}\begin{reglist}
\item [接上页]
\label{fielddesc:SPISLVRDDMADONEINTSET}\item [SPI\_SLV\_RD\_DMA\_DONE\_INT\_SET] 写 1 置位 \hyperref[int:SPISLVRDDMADONEINT]{SPI\_SLV\_RD\_DMA\_DONE\_INT} 中断。 (WT)
\label{fielddesc:SPISLVWRDMADONEINTSET}\item [SPI\_SLV\_WR\_DMA\_DONE\_INT\_SET] 写 1 置位 \hyperref[int:SPISLVWRDMADONEINT]{SPI\_SLV\_WR\_DMA\_DONE\_INT} 中断。 (WT)
\label{fielddesc:SPISLVRDBUFDONEINTSET}\item [SPI\_SLV\_RD\_BUF\_DONE\_INT\_SET] 写 1 置位 \hyperref[int:SPISLVRDBUFDONEINT]{SPI\_SLV\_RD\_BUF\_DONE\_INT} 中断。 (WT)
\label{fielddesc:SPISLVWRBUFDONEINTSET}\item [SPI\_SLV\_WR\_BUF\_DONE\_INT\_SET] 写 1 置位 \hyperref[int:SPISLVWRBUFDONEINT]{SPI\_SLV\_WR\_BUF\_DONE\_INT} 中断。 (WT)
\label{fielddesc:SPITRANSDONEINTSET}\item [SPI\_TRANS\_DONE\_INT\_SET] 写 1 置位 \hyperref[int:SPITRANSDONEINT]{SPI\_TRANS\_DONE\_INT} 中断。 (WT)
\label{fielddesc:SPIDMASEGTRANSDONEINTSET}\item [SPI\_DMA\_SEG\_TRANS\_DONE\_INT\_SET] 写 1 置位 \hyperref[int:SPIDMASEGTRANSDONEINT]{SPI\_DMA\_SEG\_TRANS\_DONE\_INT} 中断。 (WT)
\label{fielddesc:SPISEGMAGICERRINTSET}\item [SPI\_SEG\_MAGIC\_ERR\_INT\_SET] 写 1 置位 \hyperref[int:SPISEGMAGICERRINT]{SPI\_SEG\_MAGIC\_ERR\_INT} 中断。 (WT)
\label{fielddesc:SPISLVCMDERRINTSET}\item [SPI\_SLV\_CMD\_ERR\_INT\_SET] 写 1 置位 \hyperref[int:SPISLVCMDERRINT]{SPI\_SLV\_CMD\_ERR\_INT} 中断。 (WT)
\label{fielddesc:SPIMSTRXAFIFOWFULLERRINTSET}\item [SPI\_MST\_RX\_AFIFO\_WFULL\_ERR\_INT\_SET] 写 1 置位 \hyperref[int:SPIMSTRXAFIFOWFULLERRINT]{SPI\_MST\_RX\_AFIFO\_WFULL\_ERR\_INT} 中断。 (WT)
\label{fielddesc:SPIMSTTXAFIFOREMPTYERRINTSET}\item [SPI\_MST\_TX\_AFIFO\_REMPTY\_ERR\_INT\_SET] 写 1 置位 \hyperref[int:SPIMSTTXAFIFOREMPTYERRINT]{SPI\_MST\_TX\_AFIFO\_REMPTY\_ERR\_INT} 中断。 (WT)
\label{fielddesc:SPIAPP2INTSET}\item [SPI\_APP2\_INT\_SET] 写 1 置位 \hyperref[int:SPIAPP2INT]{SPI\_APP2\_INT} 中断。 (WT)
\label{fielddesc:SPIAPP1INTSET}\item [SPI\_APP1\_INT\_SET] 写 1 置位 \hyperref[int:SPIAPP1INT]{SPI\_APP1\_INT} 中断。 (WT)
\end{reglist}\end{regdesc}
\end{register}

\begin{register}{H}{SPI\_W\regindex{n}\_REG (\regindex{n}: 0-15)}{0x0098 + 0x4*\regindex{n}}\label{regdesc:SPIW0REG}
\regfield{SPI\_BUF\regindex{n}}{32}{0}{{0}}%
\reglabel{Reset}\regnewline%
\begin{regdesc}\begin{reglist}
\label{fielddesc:SPIBUF0}\item [SPI\_BUF\regindex{n}] 数据 buffer \regindex{n}，32 位。 (R/W/SS)
\end{reglist}\end{regdesc}
\end{register}


\begin{register}{H}{SPI\_DATE\_REG}{0x00F0}\label{regdesc:SPIDATEREG}
\regfield{(reserved)}{4}{28}{0000}%
\regfield{SPI\_DATE}{28}{0}{{0x2201300}}%
\reglabel{Reset}\regnewline%
\begin{regdesc}\begin{reglist}
\label{fielddesc:SPIDATE}\item [SPI\_DATE] 版本寄存器。 (R/W)
\end{reglist}\end{regdesc}
\end{register}
